
\documentclass[12pt,a4paper]{article}

% =====================================================
% PAQUETES
% =====================================================
\usepackage[utf8]{inputenc}
\usepackage[T1]{fontenc}
\usepackage[spanish,es-nodecimaldot]{babel}
\usepackage{geometry}
\usepackage{amsmath, amssymb, amsfonts}
\usepackage{booktabs}
\usepackage{longtable}
\usepackage{array}
\usepackage{xcolor}
\usepackage{hyperref}
\usepackage{enumitem}
\usepackage{listings}
\usepackage{tcolorbox}
\usepackage{fancyhdr}
\usepackage{titlesec}
\usepackage{caption}

\geometry{left=2.5cm,right=2.5cm,top=2.8cm,bottom=2.8cm}
\setlength{\parskip}{0.6em}
\setlength{\parindent}{0pt}

\hypersetup{
    colorlinks=true,
    linkcolor=black,
    urlcolor=blue,
    citecolor=black
}

% =====================================================
% COLORES Y ESTILO
% =====================================================
\definecolor{azulUP}{RGB}{0,70,127}
\definecolor{grisClaro}{RGB}{245,247,250}
\definecolor{verde}{RGB}{0,120,80}
\definecolor{morado}{RGB}{85,50,140}
\definecolor{rojo}{RGB}{160,45,45}

\titleformat{\section}{\Large\bfseries\color{azulUP}}{\thesection}{0.8em}{}[\titlerule]
\titleformat{\subsection}{\large\bfseries\color{azulUP}}{\thesubsection}{0.8em}{}
\titleformat{\subsubsection}{\normalsize\bfseries\color{azulUP}}{\thesubsubsection}{0.8em}{}

\pagestyle{fancy}
\fancyhf{}
\lhead{Calculo Estocastico}
\rhead{Sofia Escamilla}
\cfoot{\thepage}

\tcbuselibrary{skins,breakable}

\newtcolorbox{markdowncell}[1][]{
  colback=grisClaro,
  colframe=azulUP,
  title=\textbf{Celda Markdown},
  fonttitle=\bfseries\color{white},
  colbacktitle=azulUP,
  breakable,
  #1
}

\newtcolorbox{latexcell}[1][]{
  colback=white,
  colframe=morado,
  title=\textbf{Celda LaTeX / Desarrollo matematico},
  fonttitle=\bfseries\color{white},
  colbacktitle=morado,
  breakable,
  #1
}

\newtcolorbox{interpretacion}[1][]{
  colback=verde!6,
  colframe=verde,
  title=\textbf{Interpretacion financiera},
  fonttitle=\bfseries\color{white},
  colbacktitle=verde,
  breakable,
  #1
}

\newtcolorbox{alerta}[1][]{
  colback=rojo!5,
  colframe=rojo,
  title=\textbf{Nota importante},
  fonttitle=\bfseries\color{white},
  colbacktitle=rojo,
  breakable,
  #1
}

\lstdefinestyle{pythonstyle}{
    language=Python,
    basicstyle=\ttfamily\small,
    keywordstyle=\color{azulUP}\bfseries,
    commentstyle=\color{verde!70!black},
    stringstyle=\color{morado},
    showstringspaces=false,
    breaklines=true,
    frame=single,
    framerule=0.3pt,
    rulecolor=\color{azulUP!50},
    backgroundcolor=\color{grisClaro},
    tabsize=4,
    columns=fullflexible
}

\lstdefinestyle{textstyle}{
    basicstyle=\ttfamily\small,
    showstringspaces=false,
    breaklines=true,
    frame=single,
    backgroundcolor=\color{grisClaro}
}

% =====================================================
% DOCUMENTO
% =====================================================
\begin{document}

\begin{titlepage}
    \centering
    \vspace*{2cm}
    {\Large Universidad Panamericana \par}
    \vspace{0.5cm}
    {\large Licenciatura en Finanzas \par}
    \vspace{2cm}
    {\Huge\bfseries Notebook Documentado de Calculo Estocastico \par}
    \vspace{0.5cm}
    {\Large Movimiento Browniano Geometrico, Volatilidad, Monte Carlo, Opciones Americanas y Opciones Parisinas \par}
    \vspace{2cm}
    {\Large \textbf{Sofia Escamilla} \par}
    \vfill
    {\large Documento elaborado en formato \LaTeX{} con estructura de notebook para Google Colab \par}
    {\large \today \par}
\end{titlepage}

\tableofcontents
\newpage

% =====================================================
\section{Objetivo del notebook}

\begin{markdowncell}
Este documento integra los materiales de \textbf{Calculo Estocastico aplicado a Finanzas} en un solo notebook documentado. 
El objetivo es ordenar los laboratorios, ejercicios y proyecto final en una estructura compatible con Google Colab, combinando explicaciones en Markdown, desarrollo matematico en \LaTeX{} y bloques de codigo Python.

Los documentos incluidos son: Laboratorio 1 de Movimiento Browniano Geometrico, Laboratorio 2 de Volatilidad Realizada, Laboratorio 3 de Monte Carlo, Laboratorio 4 de Longstaff-Schwartz, Laboratorio 5 de Opcion Parisina, ejercicios teoricos 15.18 a 15.20 y el Proyecto 5 de Opcion Americana mediante Arbol Binomial CRR.
\end{markdowncell}

\subsection{Archivos fuente integrados}

\begin{longtable}{p{5cm}p{9cm}}
\toprule
\textbf{Archivo} & \textbf{Uso dentro del notebook} \\
\midrule
\endhead
\texttt{Laboratorio\_1\_GBM.pdf} & Estimacion de un Movimiento Browniano Geometrico con datos de AAPL. \\
\texttt{Laboratorio\_2\_Volatilidad.pdf} & Calculo de volatilidad realizada con ventanas moviles de 21, 63 y 126 dias. \\
\texttt{Laboratorio\_3\_MonteCarlo.pdf} & Valuacion Monte Carlo de una opcion call europea y comparacion con Black-Scholes. \\
\texttt{Laboratorio\_4\_LongstaffSchwartz.pdf} & Valuacion de una opcion put americana mediante Longstaff-Schwartz. \\
\texttt{Laboratorio\_5\_Opcion\_Parisina.pdf} & Simulacion de una opcion call parisina down-and-out con reloj de barrera. \\
\texttt{Ejercicios.pdf} & Ejercicios 15.18, 15.19 y 15.20 sobre put americana, opcion parisina y volatilidad historica vs implicita. \\
\texttt{Opcion\_americana.html} & Proyecto 5 sobre opcion americana mediante arbol binomial CRR. \\
\bottomrule
\end{longtable}

% =====================================================
\section{Preparacion del ambiente en Google Colab}

\begin{markdowncell}
La primera parte del notebook prepara el entorno de trabajo. Se utilizan librerias para manejo de datos, simulacion numerica, graficas, distribuciones estadisticas y, opcionalmente, descarga de informacion financiera con \texttt{yfinance}.
\end{markdowncell}

\begin{lstlisting}[style=pythonstyle]
# ==========================================
# Instalacion de librerias
# ==========================================
!pip install yfinance scipy numpy pandas matplotlib --quiet

# ==========================================
# Importacion de paquetes
# ==========================================
import numpy as np
import pandas as pd
import matplotlib.pyplot as plt
import yfinance as yf

from scipy.stats import norm
from scipy.optimize import brentq

import warnings
warnings.filterwarnings("ignore")
\end{lstlisting}

\begin{markdowncell}
Para reproducibilidad, se fija una semilla aleatoria. Esto permite que las simulaciones Monte Carlo entreguen resultados comparables cuando el notebook se vuelve a ejecutar.
\end{markdowncell}

\begin{lstlisting}[style=pythonstyle]
np.random.seed(42)
\end{lstlisting}

% =====================================================
\section{Marco teorico general}

\subsection{Movimiento Browniano}

\begin{markdowncell}
El Movimiento Browniano es el proceso estocastico base para modelar incertidumbre continua en finanzas. Su papel central es representar cambios aleatorios acumulativos en el tiempo.
\end{markdowncell}

\begin{latexcell}
Un proceso \(B_t\) es un Movimiento Browniano estandar si cumple:
\[
B_0 = 0,
\]
\[
B_t - B_s \sim N(0,t-s), \qquad 0 \leq s < t,
\]
y sus incrementos son independientes.
\end{latexcell}

\subsection{Movimiento Browniano Geometrico}

\begin{markdowncell}
El Movimiento Browniano Geometrico (GBM) se utiliza para modelar precios financieros positivos. Es la base de Black-Scholes y de muchas simulaciones de trayectorias.
\end{markdowncell}

\begin{latexcell}
La dinamica del activo bajo la medida fisica \(P\) se escribe como:
\[
dS_t = \mu S_t\,dt + \sigma S_t\,dB_t.
\]
Aplicando el Lema de Ito a \(\ln(S_t)\), se obtiene:
\[
d\ln(S_t)=\left(\mu-\frac{1}{2}\sigma^2\right)dt+\sigma dB_t.
\]
La solucion exacta discretizada es:
\[
S_{t+\Delta t}
=
S_t
\exp\left[
\left(\mu-\frac{1}{2}\sigma^2\right)\Delta t
+
\sigma \sqrt{\Delta t}Z
\right],
\qquad Z\sim N(0,1).
\]
\end{latexcell}

% =====================================================
\section{Laboratorio 1: Estimacion de un GBM}

\begin{markdowncell}
El primer laboratorio estima los parametros de deriva \(\mu\) y volatilidad \(\sigma\) de un Movimiento Browniano Geometrico usando datos historicos de AAPL. Despues se simulan trayectorias futuras bajo la solucion exacta discretizada del modelo.
\end{markdowncell}

\subsection{Estimacion de parametros}

\begin{latexcell}
Los log-retornos se calculan como:
\[
r_t = \ln\left(\frac{S_t}{S_{t-1}}\right).
\]
Si \(\hat{m}\) es la media diaria de los log-retornos y \(\hat{s}\) su desviacion estandar diaria, entonces:
\[
\hat{\sigma}=\sqrt{252}\hat{s},
\]
\[
\hat{\mu}=252\hat{m}+\frac{1}{2}\hat{\sigma}^2.
\]
\end{latexcell}

\begin{lstlisting}[style=pythonstyle]
# Descarga de precios de AAPL
ticker = "AAPL"
prices = yf.download(ticker, start="2021-01-01", end="2024-01-01", progress=False)["Close"].dropna()

# Log-retornos
log_returns = np.log(prices / prices.shift(1)).dropna()

# Estimacion de parametros anualizados
m_hat = log_returns.mean()
s_hat = log_returns.std()

sigma_hat = np.sqrt(252) * s_hat
mu_hat = 252 * m_hat + 0.5 * sigma_hat**2

print("Media diaria:", m_hat)
print("Volatilidad anual estimada:", sigma_hat)
print("Deriva anual estimada:", mu_hat)
\end{lstlisting}

\subsection{Simulacion de trayectorias}

\begin{lstlisting}[style=pythonstyle]
# Parametros de simulacion
S0 = float(prices.iloc[-1])
T = 1
n_steps = 252
n_sims = 100
dt = T / n_steps

# Matriz de trayectorias
paths = np.zeros((n_steps + 1, n_sims))
paths[0, :] = S0

for t in range(1, n_steps + 1):
    Z = np.random.normal(size=n_sims)
    paths[t, :] = paths[t-1, :] * np.exp(
        (mu_hat - 0.5 * sigma_hat**2) * dt + sigma_hat * np.sqrt(dt) * Z
    )

plt.figure(figsize=(12, 6))
plt.plot(paths, alpha=0.45)
plt.title("Simulacion futura de GBM para AAPL")
plt.xlabel("Dias")
plt.ylabel("Precio simulado")
plt.grid(True)
plt.show()
\end{lstlisting}

\begin{interpretacion}
En el laboratorio se obtuvo una deriva anual aproximada de \(0.1776\) y una volatilidad anual aproximada de \(0.2777\). La dispersion de las trayectorias aumenta conforme crece el horizonte temporal, lo cual es consistente con la acumulacion de incertidumbre en un proceso browniano.
\end{interpretacion}

% =====================================================
\section{Laboratorio 2: Volatilidad realizada y ventanas moviles}

\begin{markdowncell}
El segundo laboratorio estima la volatilidad realizada usando ventanas moviles de 21, 63 y 126 dias. La finalidad es observar como cambia el riesgo en el tiempo y detectar clustering de volatilidad.
\end{markdowncell}

\begin{latexcell}
La volatilidad realizada anualizada en una ventana de tamano \(n\) se calcula como:
\[
\sigma_t = std(r_{t-n+1},\ldots,r_t)\sqrt{252}.
\]
\end{latexcell}

\begin{lstlisting}[style=pythonstyle]
# Ventanas moviles de volatilidad realizada
vol_21 = log_returns.rolling(21).std() * np.sqrt(252)
vol_63 = log_returns.rolling(63).std() * np.sqrt(252)
vol_126 = log_returns.rolling(126).std() * np.sqrt(252)

plt.figure(figsize=(12, 6))
plt.plot(vol_21, label="21 dias")
plt.plot(vol_63, label="63 dias")
plt.plot(vol_126, label="126 dias")
plt.title("Volatilidad realizada anualizada")
plt.xlabel("Fecha")
plt.ylabel("Volatilidad anualizada")
plt.legend()
plt.grid(True)
plt.show()
\end{lstlisting}

\begin{interpretacion}
La ventana de 21 dias reacciona mas rapido ante cambios recientes, pero tambien es mas ruidosa. La ventana de 126 dias suaviza el comportamiento y muestra mejor la tendencia general del riesgo. La presencia de periodos donde la volatilidad alta se mantiene por varios dias es evidencia de clustering de volatilidad.
\end{interpretacion}

% =====================================================
\section{Laboratorio 3: Valuacion Monte Carlo de una opcion europea}

\begin{markdowncell}
El tercer laboratorio valua una opcion call europea mediante simulacion Monte Carlo bajo medida neutral al riesgo \(Q\), y compara el resultado con la formula cerrada de Black-Scholes.
\end{markdowncell}

\subsection{Dinamica neutral al riesgo}

\begin{latexcell}
Bajo la medida neutral al riesgo \(Q\), el activo sigue:
\[
dS_t = rS_t\,dt+\sigma S_t\,dB_t^Q.
\]
El precio terminal se simula como:
\[
S_T = S_0\exp\left[
\left(r-\frac{1}{2}\sigma^2\right)T+\sigma\sqrt{T}Z
\right].
\]
El precio de una call europea se obtiene como:
\[
C_0 = e^{-rT}E^Q[\max(S_T-K,0)].
\]
\end{latexcell}

\begin{lstlisting}[style=pythonstyle]
# Parametros de la opcion europea
S0 = float(prices.iloc[-1])
K = 200
r = 0.05
T = 1
sigma = float(sigma_hat)
N = 100000

Z = np.random.normal(size=N)
ST = S0 * np.exp((r - 0.5 * sigma**2) * T + sigma * np.sqrt(T) * Z)

payoffs = np.maximum(ST - K, 0)
price_mc = np.exp(-r * T) * np.mean(payoffs)
se_mc = np.std(np.exp(-r*T)*payoffs) / np.sqrt(N)

ci_low = price_mc - 1.96 * se_mc
ci_high = price_mc + 1.96 * se_mc

print("Precio Monte Carlo:", price_mc)
print("Error estandar:", se_mc)
print("IC 95%:", ci_low, ci_high)
\end{lstlisting}

\subsection{Black-Scholes}

\begin{latexcell}
La formula de Black-Scholes para una call europea es:
\[
C = S_0N(d_1)-Ke^{-rT}N(d_2),
\]
donde:
\[
d_1=
\frac{\ln(S_0/K)+(r+\frac{1}{2}\sigma^2)T}{\sigma\sqrt{T}},
\]
\[
d_2=d_1-\sigma\sqrt{T}.
\]
\end{latexcell}

\begin{lstlisting}[style=pythonstyle]
def black_scholes_call(S0, K, r, sigma, T):
    d1 = (np.log(S0 / K) + (r + 0.5 * sigma**2) * T) / (sigma * np.sqrt(T))
    d2 = d1 - sigma * np.sqrt(T)
    return S0 * norm.cdf(d1) - K * np.exp(-r*T) * norm.cdf(d2)

price_bs = black_scholes_call(S0, K, r, sigma, T)

print("Precio Black-Scholes:", price_bs)
print("Diferencia absoluta:", abs(price_mc - price_bs))
\end{lstlisting}

\begin{interpretacion}
Monte Carlo y Black-Scholes deben producir precios cercanos porque ambos se basan en la misma dinamica GBM bajo medida neutral al riesgo. La diferencia proviene del error estadistico de la simulacion.
\end{interpretacion}

% =====================================================
\section{Laboratorio 4: Opcion americana con Longstaff-Schwartz}

\begin{markdowncell}
El cuarto laboratorio aproxima el valor de una opcion put americana usando el algoritmo de Longstaff-Schwartz. La idea central es estimar el valor de continuacion mediante regresion hacia atras.
\end{markdowncell}

\subsection{Regla de ejercicio}

\begin{latexcell}
En cada fecha intermedia se compara:
\[
\text{Ejercicio inmediato}=K-S_t
\]
contra:
\[
\text{Valor de continuacion}=E^Q[e^{-r\Delta t}V_{t+\Delta t}\mid S_t].
\]
La opcion americana se valua como:
\[
V_t^{Am}=\max(K-S_t,\;C_t),
\]
donde \(C_t\) es el valor de continuacion estimado.
\end{latexcell}

\subsection{Regresion de Longstaff-Schwartz}

\begin{latexcell}
El valor de continuacion se aproxima con una base polinomial:
\[
F(S_t)=\beta_0+\beta_1S_t+\beta_2S_t^2.
\]
\end{latexcell}

\begin{lstlisting}[style=pythonstyle]
def longstaff_schwartz_put(S0, K, r, sigma, T, n_steps=252, n_sims=50000):
    dt = T / n_steps
    discount = np.exp(-r * dt)

    # Simulacion de trayectorias bajo Q
    S = np.zeros((n_steps + 1, n_sims))
    S[0, :] = S0

    for t in range(1, n_steps + 1):
        Z = np.random.normal(size=n_sims)
        S[t, :] = S[t-1, :] * np.exp((r - 0.5*sigma**2)*dt + sigma*np.sqrt(dt)*Z)

    # Payoff terminal
    cashflows = np.maximum(K - S[-1, :], 0)

    # Induccion hacia atras
    for t in range(n_steps - 1, 0, -1):
        itm = np.where(K - S[t, :] > 0)[0]

        if len(itm) > 0:
            X = S[t, itm]
            Y = cashflows[itm] * discount

            # Regresion polinomial de grado 2
            coeffs = np.polyfit(X, Y, deg=2)
            continuation = np.polyval(coeffs, X)
            exercise = K - X

            exercise_now = exercise > continuation
            ex_idx = itm[exercise_now]
            cashflows[ex_idx] = exercise[exercise_now]

        cashflows *= discount

    price = np.mean(cashflows)
    se = np.std(cashflows) / np.sqrt(n_sims)
    return price, se

price_lsm, se_lsm = longstaff_schwartz_put(S0, K, r, sigma, T)
print("Put americana Longstaff-Schwartz:", price_lsm)
print("Error estandar:", se_lsm)
\end{lstlisting}

\begin{interpretacion}
La prima por ejercicio temprano surge porque una put americana permite ejercer antes del vencimiento. Si el activo cae mucho por debajo del strike, puede ser optimo recibir \(K-S_t\) antes y reinvertir ese flujo a la tasa libre de riesgo.
\end{interpretacion}

% =====================================================
\section{Laboratorio 5: Opcion parisina y reloj de barrera}

\begin{markdowncell}
El quinto laboratorio simula una opcion call europea con barrera parisina down-and-out. La opcion se desactiva solo si el precio permanece por debajo de una barrera \(H\) durante \(D\) dias consecutivos.
\end{markdowncell}

\begin{latexcell}
La condicion de desactivacion es:
\[
S_t < H
\]
durante al menos \(D\) periodos consecutivos.
Si el activo vuelve por encima de la barrera antes de alcanzar \(D\), el reloj se reinicia.
\end{latexcell}

\begin{lstlisting}[style=pythonstyle]
def parisian_down_and_out_call(S0, K, H, r, sigma, T, D, n_steps=252, n_sims=50000):
    dt = T / n_steps
    discount = np.exp(-r*T)

    S = np.zeros((n_steps + 1, n_sims))
    S[0, :] = S0

    clock = np.zeros(n_sims)
    knocked_out = np.zeros(n_sims, dtype=bool)

    for t in range(1, n_steps + 1):
        Z = np.random.normal(size=n_sims)
        S[t, :] = S[t-1, :] * np.exp((r - 0.5*sigma**2)*dt + sigma*np.sqrt(dt)*Z)

        below = S[t, :] < H
        clock[below] += 1
        clock[~below] = 0

        knocked_out = knocked_out | (clock >= D)

    payoff = np.where(knocked_out, 0, np.maximum(S[-1, :] - K, 0))
    price = discount * np.mean(payoff)
    survival_prob = 1 - knocked_out.mean()
    knockout_prob = knocked_out.mean()

    return price, survival_prob, knockout_prob

H = 0.90 * S0
D_values = [1, 5, 10, 21, 42, 63]

results = []
for D in D_values:
    price, survival, knockout = parisian_down_and_out_call(S0, K, H, r, sigma, T, D)
    results.append([D, price, survival, knockout])

pd.DataFrame(results, columns=["D", "Precio", "Prob_Supervivencia", "Prob_Knockout"])
\end{lstlisting}

\begin{interpretacion}
Cuando \(D\) es pequeno, es mas facil que la opcion se desactive, por lo que el precio es menor. A medida que \(D\) aumenta, se vuelve mas dificil acumular suficientes dias consecutivos debajo de la barrera, por lo que la probabilidad de supervivencia aumenta y el precio de la opcion tambien.
\end{interpretacion}

% =====================================================
\section{Ejercicios 15.18 a 15.20}

\subsection{Ejercicio 15.18: Frontera de ejercicio en una put americana}

\begin{markdowncell}
Una put americana puede ejercerse antes del vencimiento. Su valor se define como el maximo entre ejercer hoy o continuar.
\end{markdowncell}

\begin{latexcell}
\[
V(t,S)=\max(K-S,\; \text{valor de continuar}).
\]
Existe una frontera critica \(S^*(t)\) tal que:
\[
S<S^*(t) \Rightarrow \text{ejercer},
\]
\[
S>S^*(t) \Rightarrow \text{continuar}.
\]
En la frontera:
\[
V(t,S^*(t))=K-S^*(t).
\]
\end{latexcell}

\begin{interpretacion}
Para una put americana, precios bajos del subyacente hacen mas atractivo ejercer porque la opcion esta muy in-the-money. Precios altos hacen mas atractivo esperar porque todavia existe valor temporal.
\end{interpretacion}

\subsection{Ejercicio 15.19: Sensibilidad de una opcion parisina}

\begin{latexcell}
Si la barrera \(H\) aumenta, la barrera queda mas cerca del precio actual:
\[
H \uparrow \Rightarrow \text{mayor probabilidad de knockout}.
\]
Por lo tanto:
\[
H \uparrow \Rightarrow \text{precio de la opcion disminuye}.
\]
Si el tiempo minimo \(D\) aumenta:
\[
D \uparrow \Rightarrow \text{knockout mas dificil}.
\]
Entonces:
\[
D \uparrow \Rightarrow \text{precio de la opcion aumenta}.
\]
\end{latexcell}

\subsection{Ejercicio 15.20: Volatilidad historica vs volatilidad implicita}

\begin{latexcell}
La volatilidad historica se estima con retornos observados:
\[
R_t=\ln\left(\frac{S_t}{S_{t-1}}\right),
\]
\[
\sigma_{hist}=\sqrt{Var(R_t)}.
\]
La volatilidad implicita se obtiene del precio de mercado de una opcion:
\[
C_{BS}(S,K,r,\sigma,T)=C_{mercado}.
\]
La \(\sigma\) que iguala el precio teorico con el precio observado es la volatilidad implicita:
\[
\sigma=\sigma_{impl}.
\]
\end{latexcell}

% =====================================================
\section{Proyecto 5: Opcion americana mediante arbol binomial CRR}

\begin{markdowncell}
El proyecto final desarrolla la valuacion de una opcion put europea y americana mediante el modelo binomial Cox-Ross-Rubinstein (CRR). La diferencia principal es que la opcion americana permite ejercicio anticipado en cada nodo.
\end{markdowncell}

\subsection{Parametros del arbol}

\begin{latexcell}
El tiempo al vencimiento \(T\) se divide en \(N\) pasos:
\[
\Delta t = \frac{T}{N}.
\]
Los factores de subida y bajada son:
\[
u=e^{\sigma\sqrt{\Delta t}},
\]
\[
d=e^{-\sigma\sqrt{\Delta t}}=\frac{1}{u}.
\]
La probabilidad neutral al riesgo es:
\[
p=\frac{e^{r\Delta t}-d}{u-d}.
\]
El factor de descuento por paso es:
\[
disc=e^{-r\Delta t}.
\]
\end{latexcell}

\subsection{Backward induction}

\begin{latexcell}
Al vencimiento, el payoff de una put es:
\[
V_{N,j}=\max(K-S_{N,j},0).
\]
Para una put europea:
\[
V_{i,j}^{Eu}=e^{-r\Delta t}\left[pV_{i+1,j+1}+(1-p)V_{i+1,j}\right].
\]
Para una put americana:
\[
V_{i,j}^{Am}=\max\left(K-S_{i,j},\; e^{-r\Delta t}\left[pV_{i+1,j+1}+(1-p)V_{i+1,j}\right]\right).
\]
\end{latexcell}

\begin{lstlisting}[style=pythonstyle]
def binomial_put_crr(S0, K, r, sigma, T, N, american=False):
    dt = T / N
    u = np.exp(sigma * np.sqrt(dt))
    d = np.exp(-sigma * np.sqrt(dt))
    p = (np.exp(r * dt) - d) / (u - d)
    disc = np.exp(-r * dt)

    # Arbol de precios
    S = np.zeros((N + 1, N + 1))
    for i in range(N + 1):
        for j in range(i + 1):
            S[i, j] = S0 * (u**j) * (d**(i-j))

    # Payoff terminal
    V = np.zeros((N + 1, N + 1))
    for j in range(N + 1):
        V[N, j] = max(K - S[N, j], 0)

    # Induccion hacia atras
    for i in range(N - 1, -1, -1):
        for j in range(i + 1):
            continuation = disc * (p * V[i+1, j+1] + (1-p) * V[i+1, j])
            if american:
                exercise = max(K - S[i, j], 0)
                V[i, j] = max(exercise, continuation)
            else:
                V[i, j] = continuation

    return V[0, 0], S, V

# Parametros del proyecto
S0_proj = 50
K_proj = 52
r_proj = 0.05
sigma_proj = 0.30
T_proj = 1
N_proj = 3

put_eu, tree_S, tree_V_eu = binomial_put_crr(S0_proj, K_proj, r_proj, sigma_proj, T_proj, N_proj, american=False)
put_am, tree_S, tree_V_am = binomial_put_crr(S0_proj, K_proj, r_proj, sigma_proj, T_proj, N_proj, american=True)

print("Put europea:", put_eu)
print("Put americana:", put_am)
print("Prima por ejercicio temprano:", put_am - put_eu)
\end{lstlisting}

\begin{interpretacion}
La put americana debe valer al menos lo mismo que la put europea:
\[
P_{Am}\geq P_{Eu}.
\]
La diferencia entre ambas se llama prima por ejercicio temprano y representa el valor de la flexibilidad de ejercer antes del vencimiento.
\end{interpretacion}

% =====================================================
\section{Resumen comparativo de laboratorios}

\begin{longtable}{p{3cm}p{5cm}p{6cm}}
\toprule
\textbf{Modulo} & \textbf{Modelo principal} & \textbf{Idea central} \\
\midrule
\endhead
Laboratorio 1 & GBM & Estimar \(\mu\) y \(\sigma\) para simular precios futuros. \\
Laboratorio 2 & Volatilidad realizada & Medir riesgo con ventanas moviles y detectar clustering. \\
Laboratorio 3 & Monte Carlo y Black-Scholes & Valuar una call europea bajo medida neutral al riesgo. \\
Laboratorio 4 & Longstaff-Schwartz & Aproximar el valor de una put americana mediante regresion hacia atras. \\
Laboratorio 5 & Opcion Parisina & Incorporar memoria temporal mediante un reloj de barrera. \\
Ejercicios & Frontera libre, barrera y volatilidad implicita & Interpretar decisiones de ejercicio y sensibilidad de derivados. \\
Proyecto 5 & Arbol binomial CRR & Valuar opcion europea y americana mediante backward induction. \\
\bottomrule
\end{longtable}

% =====================================================
\section{Conclusiones}

\begin{interpretacion}
El conjunto de laboratorios muestra una ruta completa de aplicacion de Calculo Estocastico a Finanzas. Primero se modela el precio del activo con Movimiento Browniano Geometrico y se estima su volatilidad. Despues se utiliza esa dinamica para valuar derivados mediante Monte Carlo, Black-Scholes, Longstaff-Schwartz y arboles binomiales.

La parte mas importante es entender que cada modelo responde a una necesidad distinta. GBM permite simular trayectorias; la volatilidad realizada mide riesgo observado; Monte Carlo aproxima valores esperados bajo incertidumbre; Black-Scholes da una solucion cerrada para opciones europeas; Longstaff-Schwartz aproxima decisiones de ejercicio anticipado; y las opciones parisinas agregan dependencia de trayectoria mediante memoria temporal.

En conjunto, el notebook documenta los principales metodos numericos y probabilisticos usados en Finanzas Cuantitativas para valorar derivados y analizar riesgo.
\end{interpretacion}

% =====================================================
\section{Referencias}

\begin{itemize}
    \item Black, F. y Scholes, M. (1973). \textit{The Pricing of Options and Corporate Liabilities}.
    \item Cox, J., Ross, S. y Rubinstein, M. (1979). \textit{Option Pricing: A Simplified Approach}.
    \item Hull, J. C. \textit{Options, Futures, and Other Derivatives}.
    \item Shreve, S. \textit{Stochastic Calculus for Finance II}.
    \item Longstaff, F. y Schwartz, E. (2001). \textit{Valuing American Options by Simulation}.
    \item Documentos fuente proporcionados: \texttt{Laboratorio\_1\_GBM.pdf}, \texttt{Laboratorio\_2\_Volatilidad.pdf}, \texttt{Laboratorio\_3\_MonteCarlo.pdf}, \texttt{Laboratorio\_4\_LongstaffSchwartz.pdf}, \texttt{Laboratorio\_5\_Opcion\_Parisina.pdf}, \texttt{Ejercicios.pdf}, \texttt{Opcion\_americana.html}.
\end{itemize}

% =====================================================
\appendix
\section{Apéndice: checklist para ejecutar en Colab}

\begin{enumerate}
    \item Instalar librerias: \texttt{yfinance}, \texttt{numpy}, \texttt{pandas}, \texttt{matplotlib}, \texttt{scipy}.
    \item Descargar precios de AAPL.
    \item Calcular log-retornos.
    \item Estimar \(\hat{\mu}\) y \(\hat{\sigma}\).
    \item Simular trayectorias GBM.
    \item Calcular volatilidad realizada con ventanas 21, 63 y 126 dias.
    \item Valuar una call europea con Monte Carlo.
    \item Comparar contra Black-Scholes.
    \item Valuar una put americana con Longstaff-Schwartz.
    \item Simular una opcion parisina con reloj de barrera.
    \item Construir un arbol binomial CRR.
    \item Comparar put europea vs put americana.
    \item Interpretar prima por ejercicio temprano.
\end{enumerate}

\end{document}
